%%%%%%%%%%%%%%%%%
% Fazit
%%%%%%%%%%%%%%%%%
\section{Fazit}\label{cap:conc}
Die vorangegangenen Darstellungen verdeutlichen, wie problematisch die Frage nach der empirisch besten Strategie für einen stochastischen Algorithmus ist, wenn keine hinreichend große Stichprobe generiert werden kann. 

Zudem hat sich gezeigt, dass von einer parallelen Berechnung Abstand genommen werden sollte. 
Die parallel verarbeiteten Instanzen weisen tendenziell verschlechterte Gesamtlaufzeiten auf. 
Die Brennzahlen aus \cite{NAZERI2023} konnten für jede einzelne Instanz und jede Strategie bestätigt, aber in keinem Fall unterboten werden. 
Die entscheidende Erkenntnis bleibt die Zuverlässigkeit der implementierten Lösung: 
Das Programm kann für jede der 45 untersuchten Kombinationen aus Zentralitätsmaß und Penalty-Funktion zulässige Lösungen generieren. 
Aufgrund der exzessiv langen Laufzeiten des Programms kann jedoch keine ausreichend große Stichprobe für alle 45 Strategien gewonnen werden. 
